% Source: https://artofproblemsolving.com/community/c3h1710182_junior_regional_olympiad__fbh_2017_grade_8_problem_4
% Stable identity: aops-topic-1710182-q01
\textbf{(ii) and (iii) cannot both be true.} If $n=2k+5$, then $n+k=3k+5\equiv2\pmod3$, so (iii) fails. Hence the false statement is (ii) or (iii).

\textbf{Case 1: (ii) false; (i), (iii), (iv) true.} Then $n+7k=(n+k)+6k$ is divisible by $3$ and at least $1+7=8$, so it is not prime, contradicting (iv). No solutions.

\textbf{Case 2: (iii) false; (i), (ii), (iv) true.} Now $n=2k+5$ and $k\mid n+1=2k+6$, so $k\mid6$, i.e.\ $k\in\{1,2,3,6\}$. Then $n+7k=9k+5$ equals $14,\ 23,\ 32,\ 59$ respectively, which is prime only for $k=2$ and $k=6$. For these, (iii) is indeed false by the first remark.

So $\boxed{(n,k)=(9,2)\text{ or }(17,6)}$. Check: $k=2$: $2\mid10$, $9=2\cdot2+5$, $9+14=23$ prime, $11$ not divisible by $3$; $k=6$: $6\mid18$, $17=12+5$, $17+42=59$ prime, $23$ not divisible by $3$.
